\documentclass[11pt]{article}
\usepackage[margin=1in]{geometry}
\usepackage{amsmath,amssymb}
\usepackage{booktabs}
\usepackage{natbib}
\usepackage[hidelinks]{hyperref}
% Shared notation for the differential-geometric SN Ia light-curve project.
% Input this file from every scientific document so notation stays consistent.
%
% Requires: amsmath, amssymb

% --- Number sets -------------------------------------------------------
\newcommand{\Real}{\mathbb{R}}
\newcommand{\Rn}{\mathbb{R}^{n}}
\newcommand{\Rtwo}{\mathbb{R}^{2}}
\newcommand{\Rthree}{\mathbb{R}^{3}}

% --- Differential geometry ---------------------------------------------
\newcommand{\dd}{\mathrm{d}}
\newcommand{\curve}{\gamma}                 % the light-curve curve
\newcommand{\arc}{s}                        % arclength
\newcommand{\curv}{\kappa}                  % curvature
\newcommand{\tors}{\tau}                    % torsion (n >= 3 only)
\newcommand{\frenet}[1]{\mathbf{E}_{#1}}    % Frenet frame vector
\newcommand{\tang}{\mathbf{T}}              % unit tangent
\newcommand{\norml}{\mathbf{N}}             % unit normal
\newcommand{\arclen}{L}                     % total arclength of the window

% --- Model parameters ---------------------------------------------------
\newcommand{\shape}{\theta}                 % shape latents
\newcommand{\shapeone}{\theta_{1}}
\newcommand{\shapetwo}{\theta_{2}}
\newcommand{\tscale}{\sigma}               % timing scale
\newcommand{\tmax}{t_{\mathrm{max}}}        % time origin
\newcommand{\warpone}{a_{1}}                % quadratic timing corrections
\newcommand{\warptwo}{a_{2}}
\newcommand{\dm}{\mu}                       % normalization / distance modulus
\newcommand{\colr}{c}                       % colour translation (variant only)

% --- Photometry ---------------------------------------------------------
\newcommand{\magg}{m_{g}}
\newcommand{\magr}{m_{r}}
\newcommand{\fluxg}{f_{g}}
\newcommand{\fluxr}{f_{r}}
\newcommand{\zp}{\mathrm{zp}}
\newcommand{\phase}{p}                      % rest-frame phase
\newcommand{\zred}{z}

% --- Objects and abbreviations -----------------------------------------
\newcommand{\snia}{SN\,Ia}
\newcommand{\sneia}{SNe\,Ia}
\newcommand{\ztf}{ZTF}
\newcommand{\salt}{SALT2}
\newcommand{\NN}{\mathrm{NN}}

% --- Operators ----------------------------------------------------------
\DeclareMathOperator{\SO}{SO}
\DeclareMathOperator*{\argmin}{arg\,min}
\newcommand{\deriv}[2]{\frac{\dd #1}{\dd #2}}
\newcommand{\pderiv}[2]{\frac{\partial #1}{\partial #2}}
\newcommand{\abs}[1]{\left| #1 \right|}
\newcommand{\norm}[1]{\left\| #1 \right\|}


\title{\salt\ Distilled: \\ What the Classical Model Supplies to a Geometric One}
\date{\today}

\begin{document}
\maketitle

\begin{abstract}
A working distillation of \citet{Guy2007} restricted to what bears on the
differential-geometric light-curve model. This is not a summary of the paper.
Spectroscopic recalibration, UV modelling, photometric identification and
photometric redshifts are omitted; they are central to \salt\ and irrelevant
here. What is retained is the parameterisation, the training strategy, the
error model, the distance estimator, and the specific numerical values needed
to benchmark against \salt\ on the same supernovae.
\end{abstract}

\section{The functional form}

\salt\ models the rest-frame spectral energy distribution as a linear expansion
about a mean sequence, modulated by a multiplicative colour term:
\begin{equation}
  F(\mathrm{SN}, \phase, \lambda)
    = x_{0}\left[M_{0}(\phase,\lambda) + x_{1}M_{1}(\phase,\lambda) + \dots\right]
      \times \exp\!\left[\colr\,\mathrm{CL}(\lambda)\right].
  \label{eq:salt2form}
\end{equation}
Here $\phase$ is rest-frame phase relative to $B$-band maximum and $\lambda$ is
rest-frame wavelength.

\paragraph{Global quantities, shared by all supernovae.} The mean sequence
$M_{0}(\phase,\lambda)$; the variability components $M_{k}(\phase,\lambda)$ for
$k>0$; the colour law $\mathrm{CL}(\lambda)$.

\paragraph{Per-supernova quantities.} The normalization $x_{0}$; the shape
coefficients $x_{k}$ for $k>0$; the colour $\colr$; the date of $B$-band
maximum $t_{0}$.

In the released model only two components are retained, so each supernova
carries exactly four parameters, $(x_{0}, x_{1}, \colr, t_{0})$. Additional
components were found to be poorly constrained over most of the phase space and
only marginally significant where coverage is good.

\paragraph{Colour.} The colour parameter is defined as an offset from the mean,
$\colr = (B-V)_{\mathrm{MAX}} - \langle B-V\rangle$. Critically,
\eqref{eq:salt2form} makes the colour term \emph{phase-independent}: all
remaining colour evolution with phase is carried by the linear components. The
colour law is a third-order polynomial with two free coefficients, gauge-fixed
by $\mathrm{CL}(\lambda_{B}) = 0$ and
$\mathrm{CL}(\lambda_{V}) = 0.4\log 10$.

\paragraph{Gauge on the shape parameter.} The mean and scale of $x_{1}$ are
arbitrary, since any rescaling can be absorbed into $M_{1}$. \salt\ adopts
$\langle x_{1}\rangle = 0$ and $\langle x_{1}^{2}\rangle = 1$. The geometric
model faces the identical ambiguity in its latents and adopts the same
convention.

\section{Why this is not a principal component analysis}

Equation~\eqref{eq:salt2form} is, apart from the colour exponential, a
principal component decomposition. \salt\ does not obtain it by PCA, and the
reason matters for us. PCA requires a dense, homogeneous set of observations
for every object --- in practice a spectrophotometric spectrum every four or
five days --- which does not exist. \salt\ instead fits a model with missing
data by minimising a global $\chi^{2}$ over all supernovae at once.

The geometric model inherits this constraint exactly. It cannot interpolate a
supernova's curve and then measure its curvature; it must forward-model the
sparse photometry and fit globally.

\section{Basis and resolution}

\salt\ expands the model on third-order B-splines in $(\phase,\lambda)$, chosen
for continuous second derivatives. The stated resolution is $10\times120$
parameters for $M_{0}$ (phase $\times$ wavelength) over a phase range
$[-20,+50]$~d and spectral range $[2000,9200]$~\AA, giving roughly 60~\AA\
spectral resolution, with $M_{1}$ at the coarser $10\times60$. The time axis is
remapped so resolution near maximum is about twice as fine as at the extremes
(roughly 4.5 versus 9 days).

Two lessons transfer. First, the choice of basis is irrelevant where data are
dense and controls the model only where they are sparse --- so the geometric
model must identify and report its own poorly-constrained regions rather than
trusting the network's extrapolation. Second, phase resolution should be
non-uniform, concentrated near maximum. In the geometric model this happens
naturally: arclength accumulates fastest where the curve moves fastest.

\section{Training strategy}

Training minimises a global $\chi^{2}$ over the full data set by Gauss--Newton
iteration. The procedure is explicitly \emph{staged}: the mean sequence
$M_{0}$, the colour law, and the per-supernova parameters are fitted first;
once converged, a further component is added and everything is refitted. The
authors report that convergence is insensitive to the initial guess for the
components.

This staging is directly applicable. The model ladder is not merely an
evaluation device --- it is also the training schedule. Fit L0 to convergence,
then introduce $\shapeone$, then $\shapetwo$, rather than optimising all
latents from random initialisation.

\paragraph{Regularisation.} Where a region of phase space is covered by
photometry but not spectroscopy, combinations of parameters are unconstrained.
\salt\ adds a penalty on second derivatives of each component,
\begin{equation}
  \chi^{2}_{\mathrm{REGUL}} = n \sum_{k} M_{k}^{\mathsf{T}} D^{\mathsf{T}} D M_{k},
\end{equation}
active only where data are insufficient. The normalisation $n$ was tuned on
simulated data so that the induced bias in K-corrections stays below
0.005~mag.

The method of tuning is worth copying even though the penalty itself is not:
generate synthetic data from the fitted model, retrain, and choose the
regularisation strength so the recovered bias is small compared to statistical
error. The geometric model needs the analogous procedure for its smoothness
penalty on $\curv(\arc)$ and its shrinkage on $\warpone,\warptwo$.

\section{Error model}

\salt\ treats model error as a first-class term, not an afterthought. The model
variance is a scaled version of the propagated component covariance,
\begin{align}
  V_{\mathrm{MODEL}}(x_{1},\phase,\lambda)
    &= S(\phase,\lambda)\, V_{\mathrm{MEAN}}(x_{1},\phase,\lambda), \\
  V_{\mathrm{MEAN}}(x_{1},\phase,\lambda)
    &= H_{0}^{\mathsf{T}}V_{0}H_{0}
     + x_{1}^{2}H_{1}^{\mathsf{T}}V_{1}H_{1}
     + 2x_{1}H_{0}^{\mathsf{T}}C_{0,1}H_{1},
\end{align}
where $H_{k}(\phase,\lambda)$ are the basis vectors and $S(\phase,\lambda)$ is
a scaling determined bin by bin so that the normalised residuals have unit
variance. Note the explicit dependence on $x_{1}$: model error is larger for
supernovae further from the mean.

Remaining variability beyond the components is assessed by a \textbf{jackknife}:
for each supernova, retrain on all others, then fit only that supernova's own
parameters and examine its residuals. \salt\ reports that an additional
0.022~mag must be added in quadrature to match the observed dispersion.

The jackknife is the single most important methodological import. It is
leave-one-out cross-validation, and it is the only defensible way to ask
whether an added parameter genuinely improves prediction rather than merely
absorbing noise. For a model with seven per-supernova parameters and a neural
curvature representation, in-sample residuals are meaningless.

\section{Distance estimator}

Distances follow a Tripp-style linear standardisation \citep{Tripp1998}. With
\begin{equation}
  X_{s} = \begin{pmatrix} m^{*}_{B} \\ x_{1} \\ \colr \end{pmatrix},
  \qquad
  V = \begin{pmatrix} 1 \\ \alpha \\ -\beta \end{pmatrix},
\end{equation}
the corrected distance modulus is $V^{\mathsf{T}}X_{s}$ up to an absolute
magnitude offset, i.e.
\begin{equation}
  \dm = m^{*}_{B} - M + \alpha x_{1} - \beta\,\colr .
  \label{eq:tripp2}
\end{equation}
Fitting a flat $\Lambda$CDM cosmology to the A06 sample, \salt\ obtains
\begin{equation}
  \Omega_{M} = 0.240 \pm 0.033, \qquad
  \alpha = 0.13 \pm 0.013, \qquad
  \beta = 1.77 \pm 0.16 .
\end{equation}
The variance of $m^{*}_{B}$ includes both the intrinsic dispersion and a
peculiar-velocity contribution taken at $300\ \mathrm{km\,s^{-1}}$.

Two points transfer. The geometric model's analogue of \eqref{eq:tripp2} is
$\dm_{\mathrm{corr}} = \dm - \alpha\shapeone - \beta\shapetwo + \mathcal{M}$,
and $\alpha,\beta$ are gauge-dependent: they scale inversely with whatever
normalisation is imposed on the latents, so they are comparable to \salt's
values only after adopting the matching $\langle\shape\rangle=0$,
$\langle\shape^{2}\rangle=1$ convention. Separately, the peculiar-velocity term
enters only at the Hubble-diagram stage; it does not belong in the light-curve
likelihood, where $\dm$ is free per supernova and absorbs it exactly.

\section{The stretch correspondence}

\salt\ reports that most of the variability captured by $x_{1}$ amounts to a
simple stretching of the light curves, despite no such behaviour being imposed.
The conversions given are
\begin{align}
  s^{(\mathrm{SALT})} &= 0.98 + 0.091x_{1} + 0.003x_{1}^{2} - 0.00075x_{1}^{3}, \\
  s^{(\mathrm{G01})}  &= 1.07 + 0.069x_{1} - 0.015x_{1}^{2} + 0.00067x_{1}^{3}, \\
  \Delta m_{15}       &= 1.09 - 0.161x_{1} + 0.013x_{1}^{2} - 0.00130x_{1}^{3}.
\end{align}

This is the most consequential result in the paper for our purposes, and it
cuts both ways. It is encouraging, in that a model given free rein to find
components recovers something close to a pure time reparameterisation --- which
is precisely the degree of freedom the geometric model promotes to a
first-class object in $\arc(t)$. It is also a warning: if $x_{1}$ is
substantially a stretch, then a model that separates timing from shape should
find that $\shapeone$ is largely \emph{redundant} with $\tscale$. A strong
fitted correlation between them would confirm the interpretation; the absence
of one would demand explanation.

These conversions also supply the cleanest interpretability test available.
Correlating fitted $(\tscale, \shapeone, \shapetwo)$ against \salt's
$(x_{1}, \colr)$ and against $\Delta m_{15}$ tells us what the geometric
latents have actually learned.

\section{What is deliberately not adopted}

\begin{center}
\begin{tabular}{ll}
\toprule
\salt\ element & Status here \\
\midrule
SED surface over $(\phase,\lambda)$ & replaced by a curve in band space \\
B-spline basis                      & replaced by a neural representation \\
Spectroscopic training data         & not used; photometry only \\
Spectral recalibration polynomials  & not applicable \\
Colour law $\mathrm{CL}(\lambda)$   & not identifiable with two bands \\
K-corrections from the SED          & avoided by the $\zred<0.05$ cut \\
Staged component training           & \textbf{adopted} \\
Jackknife validation                & \textbf{adopted} \\
Latent gauge $\langle x\rangle=0$, $\langle x^{2}\rangle=1$ & \textbf{adopted} \\
Tripp standardisation               & \textbf{adopted} \\
\bottomrule
\end{tabular}
\end{center}

The colour law entry deserves comment. With only $g$ and $r$, there is no
wavelength axis along which to fit $\mathrm{CL}(\lambda)$; colour is a single
number, not a law. This is why the geometric model must decide by experiment
whether colour needs its own translational degree of freedom or can be absorbed
into curvature, rather than importing a functional form from \salt.

\section{Role of \salt\ in this project}

\salt\ is used strictly \textbf{downstream}. Fits are performed independently
on the same supernovae and compared against the geometric model on held-out
Hubble residuals. \salt\ is never used to preprocess, K-correct, interpolate,
or initialise. Were it used upstream, its SED errors would propagate silently
into the data and the claim of an independent alternative would fail.

\bibliographystyle{plainnat}
\bibliography{refs}

\end{document}
